% 04 Apr: AI clean

\section{Augsburg and the Minority}

"Folkebladet," July 25, 1894

"... After a brief consultation with some of Augsburg’s supporters, who gathered in the hall that evening, and in which it was unanimously declared that, in whatever was yielded, Augsburg’s principle regarding congregational life and pastoral training must not be relinquished, I traveled to St. Paul the following day, praying that God would guide matters so that they might serve his church for blessing and growth. I was determined to sacrifice everything, even myself, in order to bring about understanding and sincere cooperation between the majority and the minority—but principles could not be given up, nor the cause to which we have committed ourselves for these twenty years.

..........

I was present during the annual meeting’s consideration of this matter. Not only the voting, but still more the statements made on that occasion, convinced me that there exists—notwithstanding the personal bitternesses and plans of revenge introduced by the old party—a deep gulf of principle between the majority and the minority, or Augsburg’s supporters, precisely in the question of pastoral training and the inseparable questions of the ministry and the congregation, a gulf which I now clearly see cannot be smoothed out and overcome except by persistent and patient labor, in order by way of development, with God’s help, to make the free-congregation idea work like leaven within the majority.

That the minority’s two theological professors would be forced out if the majority gained control may surely be regarded as self-evident. Of the three theological professors at the majority’s school, Prof. Schmidt is opposed to the Augsburg principle; Prof. Lund scarcely seems to have any principle in this matter except that of the General Council, and that is very far from Augsburg’s. Thus only Bøckman remains, who, even if he had been in agreement with Augsburg’s educational principle in more than a merely formal respect, was plainly not in a position, with the two men named at his side, to carry through that principle of pastoral training, and that understanding of the congregation, which has hitherto been the prevailing one at Augsburg."

Great changes must surely take place in the majority before there can again be any question of negotiations for transfer. There will first have to be serious, conscientious, perhaps hard-fought negotiations about the principle, so that the same error which Prof. Bøckman admitted the Conference had negligently committed in Scandinavia may not be repeated.

In other words, a more solid, more dependable, and above all more spiritual, more unequivocal basis in keeping with the free-congregation principle will be required if the majority and the minority are to come into true organic union and cooperation with one another.

It is the principles that must now be tried and fought through, until the truth is vindicated in our American Free Church.

Sven Oftedal
